\section{Limitations}
\label{sec:limitations}

First, there is no ground truth for ``owned'' values. The method measures response posture, not inner authenticity. This is a deliberate ceiling on the claim.

Second, clamping is ambiguous. A clamped result may mean that a model reliably returns to the assistant-service surface under this instrument, or that our role-negation prompt is too weak to cross its surface. Stronger or different cache-breaking mechanisms are future work.

Third, the coding is model-assisted. Each layer uses three model coders and consensus rules, but the coders are themselves language models. Human spot-coding, especially prompt-masked coding, would strengthen the validation.

Fourth, the world-change prompt is not equivalent to direct stated-values elicitation. It elicits normative wishes. This is values-revealing, but it should not be collapsed into ``what the model says it cares about.''

Fifth, release-date trend analysis is exploratory. Model versions, endpoint behaviour, and route/provider differences can shift over time. Lab-level trend charts are descriptive aids, not causal evidence about training trajectories.
